\chapter{Discussion and Evaluation of Results}
\label{cha:results}

This chapter turns the twelve experimental suites from
Chapter~\ref{cha:exper} into evidence for the central claim of this
report: the best-so-far trajectory reveals useful structure even when
an algorithm does not reach the optimum within the available budget.
Instead of asking only whether a run succeeds, the discussion asks how
far the best solution found so far remains from the optimum, how quickly
that gap closes, and where progress stalls.

The main signal used for this purpose is incumbent cumulative regret
(\acrshort{cr}). A rising \acrshort{cr} curve means the current
best-so-far solution is still leaving fitness on the table; a curve that
flattens means the remaining gap is closing. This makes the plots useful
in cases where success probabilities are near zero and
time-to-first-optimum summaries have little to say.

The results are organized into five thematic groups. The chapter begins
with baseline comparisons on standard pseudo-Boolean benchmarks
(Sections~\ref{sec:results:baseline}--\ref{sec:results:suite02}), then
checks the theoretical connection between runtime profiles and regret
(Section~\ref{sec:results:verification}). It then studies barrier
structure and mutation rate (Section~\ref{sec:results:barrier}), NK
ruggedness and simulated-annealing cooling schedules
(Section~\ref{sec:results:nk}), and finally hierarchical structure in
\acrshort{hiff} (Section~\ref{sec:results:hierarchical}). Where runtime
profile plots are available, they are used as a second view of the same
behavior rather than as a separate metric.


\section{Baseline Suites}
\label{sec:results:baseline}
A central diagnostic in the baseline suites is the slope of the
cumulative regret curve. Let $\rho_t^{\mathrm{inc}} = f^{*} -
f(\hat{x}_t)$ denote the incumbent instantaneous regret after
evaluation $t$, where $\hat{x}_t = \operatorname{arg\,max}_{s \leq t}
f(x_s)$ is the best solution found so far, and let
$\mathrm{CR}(T) = \sum_{t=1}^{T} \rho_t^{\mathrm{inc}}$. Since
$\mathrm{CR}(T) - \mathrm{CR}(T-1) = \rho_T^{\mathrm{inc}}$, a
persistent positive slope in the linear \acrshort{cr} plot is
direct evidence that the incumbent remains suboptimal
(Figures~\ref{fig:baseline-cr-jump3} and~\ref{fig:baseline-cr-trap10}). When the
best-so-far solution has stagnated at a local optimum, on a plateau,
or in a deceptive basin, the slope at time $T$ estimates the residual
fitness gap $f^{*} - f(\hat{x}_T)$ to the true optimum.

Conversely, if $\mathrm{CR}(T)$ approaches a horizontal asymptote
$C = \lim_{T \to \infty} \mathrm{CR}(T)$ as $T \to \infty$, then
$\rho_T^{\mathrm{inc}} \to 0$, which on the finite benchmark
landscapes studied here corresponds to eventual discovery of the
global optimum (Figures~\ref{fig:baseline-cr-onemax},
\ref{fig:baseline-cr-leadingones}, and~\ref{fig:baseline-cr-twomax}).
In finite-budget experiments the
asymptote can only be
approximated, but the approximation improves as the budget extends
beyond the typical hitting time. In the infinite-budget limit described by
Corollary~\ref{cor:asymptotic-regret},
$C = \sum_{v=1}^{f^{*}} \mathbb{E}[\tau_v] - f^{*}$, so the
asymptote is determined by the sum of expected first-hitting times
across \emph{all} intermediate fitness levels, not only the terminal
one. It should therefore be read as the total cumulative cost paid
before convergence: the area under the entire incumbent \acrshort{ir}
trajectory (Figure~\ref{fig:baseline-cr-binval}). This distinguishes
it from time-to-first-optimum
(\acrshort{ttfo}), which captures only the moment of final arrival.

All experiments in this section use $n = 100$, except for
\textsc{BinVal} at $n = 50$\footnote{The restriction arises because
  \textsc{BinVal}'s optimum is $2^n - 1$: at $n = 100$ this value
  exceeds the precision range of IEEE~754 double-precision
  floating-point arithmetic, introducing rounding error into the regret
  computation. At $n = 50$ the optimum $2^{50} - 1 \approx 1.13 \times
10^{15}$ is representable exactly.}.

\begin{figure}[!htbp]
  \centering
  \begin{subfigure}[t]{0.48\textwidth}
    \includegraphics[width=\linewidth]{results/baseline/base/onemax_history_regret_cumulative_best_1.pdf}
    \caption{\textsc{OneMax}.}
    \label{fig:baseline-cr-onemax}
  \end{subfigure}
  \begin{subfigure}[t]{0.48\textwidth}
    \includegraphics[width=\linewidth]{results/baseline/base/twomax_history_regret_cumulative_best_1.pdf}
    \caption{\textsc{TwoMax}.}
    \label{fig:baseline-cr-twomax}
  \end{subfigure}
  \begin{subfigure}[t]{0.48\textwidth}
    \includegraphics[width=\linewidth]{results/baseline/base/leadingones_history_regret_cumulative_best_1.pdf}
    \caption{\textsc{LeadingOnes}.}
    \label{fig:baseline-cr-leadingones}
  \end{subfigure}
  \begin{subfigure}[t]{0.48\textwidth}
    \includegraphics[width=\linewidth]{results/baseline/base/binval_history_regret_cumulative_best_1.pdf}
    \caption{\textsc{BinVal}.}
    \label{fig:baseline-cr-binval}
  \end{subfigure}
  \begin{subfigure}[t]{0.48\textwidth}
    \includegraphics[width=\linewidth]{results/baseline/base/jump-3_history_regret_cumulative_best_1.pdf}
    \caption{\textsc{Jump}-$3$.}
    \label{fig:baseline-cr-jump3}
  \end{subfigure}
  \begin{subfigure}[t]{0.48\textwidth}
    \includegraphics[width=\linewidth]{results/baseline/base/trap-10_history_regret_cumulative_best_1.pdf}
    \caption{\textsc{Trap}-$10$.}
    \label{fig:baseline-cr-trap10}
  \end{subfigure}
  \caption[Suite 01 \acrshort{cr} trajectories]{Suite~01
    incumbent \acrshort{cr} trajectories for the
    baseline pseudo-Boolean problems. Smooth landscapes approach a
    horizontal floor, while barrier and deceptive landscapes retain a
  positive slope under the displayed budget.}
  \label{fig:baseline-cr}
\end{figure}

\subsection{Suite 01: Algorithm-Family Comparison}
\label{sec:results:suite01}
For the simpler \textsc{OneMax} and \textsc{TwoMax} instances, the
relative ordering of mean algorithm performance aligns with the
horizontal convergence points of the median incumbent \acrshort{cr}
curves (Figures~\ref{fig:baseline-cr-onemax}
and~\ref{fig:baseline-cr-twomax}). However, on
\textsc{LeadingOnes} and \textsc{BinVal}, a structural inversion
emerges (Figures~\ref{fig:baseline-cr-leadingones}
and~\ref{fig:baseline-cr-binval}): the algorithms reach the
global optimum in the order \acrshort{rls}, \acrshort{sa}-Log
(d=2.0), \acrshort{rlse}, then $(1+1)$-\acrshort{ea}, yet their
\acrshort{cr} asymptotes are ordered in reverse.

This inversion is consistent with the mechanism described by
Corollary~\ref{cor:asymptotic-regret}:
the asymptote $C = \lim_{T \to \infty} \mathbb{E}[\mathrm{CR}(T)] =
\sum_{v=1}^{f^{*}} \mathbb{E}[\tau_v] - f^{*}$ is governed by
first-hitting times across \emph{all} intermediate fitness
levels. \acrshort{rls}'s greedy trajectory
reaches $f^{*}$ efficiently but lingers at lower fitness levels during
its ascent, inflating the per-level $\mathbb{E}[\tau_v]$ terms and
therefore $C$. The $(1+1)$-\acrshort{ea}'s multi-bit mutation
occasionally promotes the incumbent earlier in the run, reducing those
intermediate contributions and shrinking $C$, despite the later arrival
at $f^{*}$.

For \textsc{Jump} ($k=3$) and \textsc{Trap} ($k=10$), all algorithms
remain trapped in suboptimal basins within the given budget, so
$P(\mathrm{opt}) \approx 0$ and \acrshort{ttfo} is undefined for the
vast majority of runs. The incumbent \acrshort{cr} curves nevertheless
retain diagnostic structure. Each curve grows at an approximately
constant rate, confirming stagnation of the incumbent; the slope at
time $T$ estimates the residual fitness gap $\rho_T^{\mathrm{inc}} =
f^{*} - f(\hat{x}_T)$, so a steeper slope indicates that the
algorithm is trapped further from the optimum and is therefore
incurring a greater per-evaluation cost. The slopes differ across
algorithms, reflecting the fact that each stagnates at a different
local optimum depending on its search strategy, while remaining
non-zero throughout the budget because the landscape barrier (whose
depth is a fixed property of the problem) prevents any
algorithm from closing the gap. This illustrates a key advantage of
\acrshort{cr} over \acrshort{ttfo}: even when convergence is never
observed, the regret trajectory continues to distinguish algorithms
by the quality of their trapped incumbent, information that a binary
optimum achievement measure cannot provide.



\subsection{Suite 02: Extended Pseudo-Boolean Benchmark}
\label{sec:results:suite02}
%
Suite~02 shows that the Suite~01 ordering generalizes only on the smooth
part of the benchmark.

\begin{figure}[!htbp]
  \centering
  \begin{subfigure}[t]{0.48\textwidth}
    \includegraphics[width=\linewidth]{results/baseline/ext/onemax_convergence_probability.pdf}
    \caption{\textsc{OneMax} convergence probability.}
    \label{fig:suite02-onemax-cdp}
  \end{subfigure}
  \begin{subfigure}[t]{0.48\textwidth}
    \includegraphics[width=\linewidth]{results/baseline/ext/onemax_comparison_heatmap.pdf}
    \caption{\textsc{OneMax} mean-regret heat map.}
    \label{fig:suite02-onemax-heatmap}
  \end{subfigure}
  \begin{subfigure}[t]{0.48\textwidth}
    \includegraphics[width=\linewidth]{results/baseline/ext/leadingones_convergence_probability.pdf}
    \caption{\textsc{LeadingOnes} convergence probability.}
    \label{fig:suite02-leadingones-cdp}
  \end{subfigure}
  \begin{subfigure}[t]{0.48\textwidth}
    \includegraphics[width=\linewidth]{results/baseline/ext/leadingones_comparison_heatmap.pdf}
    \caption{\textsc{LeadingOnes} mean-regret heat map.}
    \label{fig:suite02-leadingones-heatmap}
  \end{subfigure}
  \begin{subfigure}[t]{0.48\textwidth}
    \includegraphics[width=\linewidth]{results/baseline/ext/twomax_convergence_probability.pdf}
    \caption{\textsc{TwoMax} convergence probability.}
    \label{fig:suite02-twomax-cdp}
  \end{subfigure}
  \begin{subfigure}[t]{0.48\textwidth}
    \includegraphics[width=\linewidth]{results/baseline/ext/twomax_comparison_heatmap.pdf}
    \caption{\textsc{TwoMax} mean-regret heat map.}
    \label{fig:suite02-twomax-heatmap}
  \end{subfigure}
  \caption[Suite 02 easy-tier summary plots]{Suite~02 easy-tier
    summary plots. Convergence probabilities and
    mean-regret heat maps show that the baseline algorithm ordering is stable
  on smooth benchmark problems under the extended budget.}
  \label{fig:suite02-easy-summary}
\end{figure}

\begin{figure}[!htbp]
  \centering
  \begin{subfigure}[t]{0.48\textwidth}
    \includegraphics[width=\linewidth]{results/baseline/ext/jump-5_history_regret_cumulative_best.pdf}
    \caption{\textsc{Jump}-$5$ cumulative regret.}
    \label{fig:suite02-jump5-cr}
  \end{subfigure}
  \begin{subfigure}[t]{0.48\textwidth}
    \includegraphics[width=\linewidth]{results/baseline/ext/trap-5_history_regret_cumulative_best.pdf}
    \caption{\textsc{Trap}-$5$ cumulative regret.}
    \label{fig:suite02-trap5-cr}
  \end{subfigure}
  \begin{subfigure}[t]{0.80\textwidth}
    \includegraphics[width=\linewidth]{results/baseline/ext/petersen-3-color_history_regret_cumulative_best.pdf}
    \caption{\textsc{Petersen}-3-colour cumulative regret.}
    \label{fig:suite02-petersen-cr}
  \end{subfigure}
  \caption[Suite 02 harder-instance regret trajectories]{Selected
    Suite~02 \acrshort{cr} trajectories for harder
    benchmark instances. The barrier problems retain sustained regret growth,
  while the graph-colouring instance is solved within the displayed budget.}
  \label{fig:suite02-selected-cr}
\end{figure}

On \textsc{OneMax}, \textsc{LeadingOnes}, and
\textsc{TwoMax}, the same separation between the four algorithm families is
visible in both convergence probability and the mean-regret heat map
Figures~\ref{fig:suite02-easy-summary}, so the
baseline ranking is stable under the extended budget.

The barrier problems change the picture.
Persistent \acrshort{rlse} exploration with $\varepsilon=0.05$
improves barrier escape.
On \textsc{Jump-3}, this changes the post-stagnation ordering of the
incumbent \acrshort{cr} curves: with fixed $\varepsilon=0.05$,
\acrshort{rlse} receives more frequent opportunities to leave the local
basin than under the Suite~01 default $\varepsilon=0.01$.

The ``hotter'' SA schedule ($d=6.0$) on \textsc{Jump-5} and \textsc{Trap-5}
reduces premature freezing, as seen by higher \acrshort{iqr} in
Figures~\ref{fig:suite02-jump5-cr} and~\ref{fig:suite02-trap5-cr}.

On \acrshort{hiff}-64 no method reaches the optimum within budget, so this
instance is retained as a hard hierarchical case rather than a convergence
benchmark.

In contrast, \textsc{Petersen-3-Color} is solved by all four
algorithms within budget, with convergence concentrated around
200--300 evaluations (Figure~\ref{fig:suite02-petersen-cr}).

\section{Verification Suite}
\label{sec:results:verification}
The verification compares two curves indexed by evaluation time: the
empirical mean cumulative regret $\bar R_t$ and the profile-derived
quantity $\hat R_t$. Since Theorem~\ref{thm:regret-profile} asserts an
identity over the full trajectory, not just at the final horizon, the
comparison is performed pointwise over $t=1,\ldots,T$.

Three complementary error summaries are reported
(Table~\ref{tab:vermeas}). The \acrfull{rmse} measures the
typical global discrepancy between the two curves and is sensitive to
persistent deviations across the horizon. The
\acrfull{mae}~\cite{higham2002acc} records the worst pointwise
discrepancy and therefore checks
whether the identity fails locally at any evaluation index. Finally, the
\acrfull{rmae} normalizes this worst-case gap
by the scale of the empirical cumulative regret curve, allowing results
to be compared across problems with very different regret magnitudes.
Together, these measures test both average agreement and worst-case
pointwise agreement.

\begin{table}[!htbp]
  \centering
  \begin{tabular}{lll}
    \toprule
    Measure & Formula & Interpretation \\
    \midrule
    \acrshort{rmse}
    & $\sqrt{\dfrac{1}{T}\sum_{t=1}^{T}\left(\hat{R}_t-\bar{R}_t\right)^2}$
    & Average curve-level discrepancy \\[1.4em]

    \acrshort{mae}
    & $\max_t \left|\hat{R}_t-\bar{R}_t\right|$
    & Worst pointwise discrepancy \\[0.8em]

    \acrshort{rmae}
    & $\dfrac{\max_t \left|\hat{R}_t-\bar{R}_t\right|}{\max_t \bar{R}_t}$
    & Scale-normalized worst-case discrepancy \\
    \bottomrule
  \end{tabular}
  \caption[Profile verification error measures]{Error measures used
    to compare profile-derived ($\hat{R}_t$) and
    empirical ($\bar{R}_t$) expected cumulative regret
    $\mathbb{E}[\mathrm{CR}(T)]$
  at each evaluation index $t$.}
  \label{tab:vermeas}
\end{table}

The verification is accepted if the profile-derived and empirical curves
satisfy two conditions: first, the relative maximum pointwise error must
remain below $1\%$,
\[
  \frac{\max_t |\hat R_t-\bar R_t|}{\max_t \bar R_t} < 0.01,
\]
and second, the normalized RMSE,
\[
  \frac{\mathrm{RMSE}}{\max_t \bar R_t},
\]
must be of the same order or smaller. These thresholds are deliberately
scale-normalized, so that the criterion is not biased against problems
such as \textsc{LeadingOnes}, where cumulative regret is naturally much
larger in absolute value.

\subsection{Suite 08: Runtime Profile Verification}
\label{sec:results:suite08}

\begin{figure}[!htbp]
  \centering
  \begin{subfigure}[t]{0.80\textwidth}
    \includegraphics[width=\linewidth]{results/verification/cr_prof_ver_onemax.pdf}
    \caption[\textsc{OneMax} verification plots]{Verification plots
      for \acrshort{rls} (black) and
      $(1+1)$-\acrshort{ea} (blue) for the \textsc{OneMax} problem,
    with maximum budget 2\,000}
  \end{subfigure}
  \begin{subfigure}[t]{0.80\textwidth}
    \includegraphics[width=\linewidth]{results/verification/cr_prof_ver_leadingones.pdf}
    \caption[\textsc{LeadingOnes} verification plots]{Verification
      plots for \acrshort{rls} (black) and
      $(1+1)$-\acrshort{ea} (blue) for the \textsc{LeadingOnes}
    problem, with maximum budget 10\,000}
  \end{subfigure}
  \caption[Empirical and profile-derived verification curves]{The
    verification curves for the empirically calculated
    mean \acrshort{cr} ($\mathbb{E}[\mathrm{CR}(T)]$, solid plot) and
  the profile-derived expected \acrshort{cr} ($\hat{R}_t$, dotted plot)}
\end{figure}

\begin{table}[!htbp]
  \centering
  \small
  \begin{tabular}{llccc}
    \toprule
    Problem & Algorithm & \acrshort{rmse} & \acrshort{mae} & \acrshort{rmae} \\
    \midrule
    \multirow{2}{8.5em}{\textsc{OneMax} (Budget: 2\,000)}
      & $(1+1)$-\acrshort{ea} & $8.90\times10^{-12}$ & $1.18\times10^{-11}$ & $2.09\times10^{-15}$ \\
      & \acrshort{rls}        & $3.02\times10^{-12}$ & $4.55\times10^{-12}$ & $9.33\times10^{-16}$ \\
    \midrule
    \multirow{2}{8.5em}{\textsc{LeadingOnes} (Budget: 10\,000)}
      & $(1+1)$-\acrshort{ea} & $1.37\times10^{-9}$ & $2.82\times10^{-9}$ & $1.22\times10^{-14}$ \\
      & \acrshort{rls}        & $3.57\times10^{-9}$ & $5.76\times10^{-9}$ & $2.33\times10^{-14}$ \\
    \bottomrule
  \end{tabular}
  \caption[Pointwise profile verification distances]{Pointwise
    distances between the empirically calculated $\mathbb{E}[\mathrm{CR}]$
    and the distributional profile-derived $\mathrm{CR}$.}
  \label{tab:profile-verification-distances}
\end{table}
\FloatBarrier

The profile identity holds to floating-point precision across all four
combinations. The largest \acrshort{mae} is $5.76\times10^{-9}$, and the
largest \acrshort{rmae} is only $2.33\times10^{-14}$.
Thus, on integer-valued benchmarks with exact fitness levels and complete time
grids, the direct trajectory average and the profile-derived tail-sum curve are
numerically indistinguishable (Table~\ref{tab:profile-verification-distances}).

This agreement is stronger than what would be expected from an independent
Monte-Carlo comparison because the two curves are computed from the same
$R=100$ runs. Sampling variation affects which trajectories are observed, but
both curves are deterministic transformations of that same sample. The
pointwise comparison therefore tests implementation agreement: for OneMax and
LeadingOnes, the profile calculation uses the exact integer fitness levels
$1,\ldots,100$ and every evaluation time point at the displayed budgets. The
remaining differences are therefore numerical round-off. 
The separate discretization caveat applies later to
real-valued landscapes such as NK.
\begin{figure}[!htbp]
  \centering
  \begin{subfigure}[t]{0.49\textwidth}
    \includegraphics[width=\linewidth]{results/verification/inv_prof_surf_onemax_RLS.pdf}
    \caption[\acrshort{rls} profile on
    \textsc{OneMax}]{\acrshort{rls} for the \textsc{OneMax} problem, with
    displayed profile horizon 1\,000}
  \end{subfigure}
  \begin{subfigure}[t]{0.49\textwidth}
    \includegraphics[width=\linewidth]{results/verification/inv_prof_surf_onemax_(1+1)-EA.pdf}
    \caption[$(1+1)$-\acrshort{ea} profile on
    \textsc{OneMax}]{$(1+1)$-\acrshort{ea} for the \textsc{OneMax} problem,
    with displayed profile horizon 1\,000}
  \end{subfigure}
  \begin{subfigure}[t]{0.48\textwidth}
    \includegraphics[width=\linewidth]{results/verification/inv_prof_surf_leadingones_RLS.pdf}
    \caption[\acrshort{rls} profile on
    \textsc{LeadingOnes}]{\acrshort{rls} for the \textsc{LeadingOnes} problem,
    with maximum budget 10\,000}
  \end{subfigure}
  \begin{subfigure}[t]{0.48\textwidth}
    \includegraphics[width=\linewidth]{results/verification/inv_prof_surf_leadingones_(1+1)-EA.pdf}
    \caption[$(1+1)$-\acrshort{ea} profile on
    \textsc{LeadingOnes}]{$(1+1)$-\acrshort{ea} for the \textsc{LeadingOnes}
    problem, with maximum budget 10\,000}
  \end{subfigure}
  \caption[Inverse runtime profile plots]{Inverse runtime profile
    [$P(\tau_v \le T)$] plots for
  \acrshort{rls} and $(1+1)$-\acrshort{ea}}
  \label{fig:invprof}
\end{figure}
The inverse runtime profile surfaces in Figure~\ref{fig:invprof} show
the expected separation between
the two benchmark regimes. For \textsc{OneMax}, the saturation frontier is
compressed into the early part of the budget: low and intermediate fitness
levels satisfy $P(\tau_v \le T) \approx 1$ after only a few hundred
evaluations, and the contour near the optimum $f^\ast$ reaches high
probability well before the end of the displayed $T=1000$ window for both
\acrshort{rls} and the $(1+1)$-\acrshort{ea}. In contrast, the
\textsc{LeadingOnes} surfaces exhibit a broad diagonal transition band
across the full $T=10000$ range. The probability mass moves gradually from
low to high fitness levels, and the saturation frontier for $v=f^\ast$
appears only in the later part of the budget. This delayed frontier is
consistent with the sequential structure of \textsc{LeadingOnes}, where
progress depends on extending the correct prefix, and with the known
asymptotic contrast between the $\Theta(n \log n)$ optimization time of
\textsc{OneMax} and the $\Theta(n^2)$ behavior of
\textsc{LeadingOnes}. Thus, the inverse runtime profiles verify the profile
construction itself and preserve the qualitative finite-$n$ runtime ordering of
the underlying benchmark problems.

The verification result supports the use of profile-derived
$\mathbb{E}[\mathrm{CR}(T)]$ curves on integer-valued or explicitly
discretized fitness levels in the remaining evaluation suites. The same profile
construction provides a useful diagnostic reference in Suites~03 and~12, where
direct averaging is more susceptible to sampling noise because only $R=25$ runs
are available. For the real-valued NK landscapes, however, those profile
surfaces should be read as discretized approximations rather than exact
identities. Section~\ref{sec:results:synthesis} returns to this point by
relating the Suite~08 verification to the saturation frontiers observed in the
NK benchmark suites.
%
\section{Barrier Structure and Mutation Rate}
\label{sec:results:barrier}
\subsection{Suite 04: Trap --- Deception Width and Mutation Rate}
\label{sec:results:suite04}
\begin{figure}[!htbp]
  \centering
  \begin{subfigure}[t]{0.48\textwidth}
    \includegraphics[width=\linewidth]{results/barrier_mut/trap/k5_history_regret_cumulative_best.pdf}
    \caption{\textsc{Trap}-$5$.}
    \label{fig:suite04-trap-k5-cr}
  \end{subfigure}
  \begin{subfigure}[t]{0.48\textwidth}
    \includegraphics[width=\linewidth]{results/barrier_mut/trap/k8_history_regret_cumulative_best.pdf}
    \caption{\textsc{Trap}-$8$.}
    \label{fig:suite04-trap-k8-cr}
  \end{subfigure}
  \begin{subfigure}[t]{0.48\textwidth}
    \includegraphics[width=\linewidth]{results/barrier_mut/trap/k10_history_regret_cumulative_best.pdf}
    \caption{\textsc{Trap}-$10$.}
    \label{fig:suite04-trap-k10-cr}
  \end{subfigure}
  \begin{subfigure}[t]{0.48\textwidth}
    \includegraphics[width=\linewidth]{results/barrier_mut/trap/k20_history_regret_cumulative_best.pdf}
    \caption{\textsc{Trap}-$20$.}
    \label{fig:suite04-trap-k20-cr}
  \end{subfigure}
  \caption[Suite 04 Trap \acrshort{cr} trajectories]{Suite~04
    incumbent \acrshort{cr} trajectories on
    \textsc{Trap} as deception width increases. The trajectories remain
    non-convergent within budget, so the ramp height and slope carry the main
  diagnostic information.}
  \label{fig:suite04-trap-cr}
\end{figure}

\begin{figure}[!htbp]
  \centering
  \begin{subfigure}[t]{0.48\textwidth}
    \includegraphics[width=\linewidth]{results/barrier_mut/trap/k5_regret_boxplot_b10000.pdf}
    \caption{\textsc{Trap}-$5$.}
    \label{fig:suite04-trap-k5-box}
  \end{subfigure}
  \begin{subfigure}[t]{0.48\textwidth}
    \includegraphics[width=\linewidth]{results/barrier_mut/trap/k8_regret_boxplot_b10000.pdf}
    \caption{\textsc{Trap}-$8$.}
    \label{fig:suite04-trap-k8-box}
  \end{subfigure}
  \begin{subfigure}[t]{0.48\textwidth}
    \includegraphics[width=\linewidth]{results/barrier_mut/trap/k10_regret_boxplot_b10000.pdf}
    \caption{\textsc{Trap}-$10$.}
    \label{fig:suite04-trap-k10-box}
  \end{subfigure}
  \begin{subfigure}[t]{0.48\textwidth}
    \includegraphics[width=\linewidth]{results/barrier_mut/trap/k20_regret_boxplot_b10000.pdf}
    \caption{\textsc{Trap}-$20$.}
    \label{fig:suite04-trap-k20-box}
  \end{subfigure}
  \caption[Suite 04 Trap regret distributions]{Suite~04 regret
    distributions at budget $10{,}000$ for
    increasing \textsc{Trap} widths. The boxplots summarize the end-of-budget
    regret separation between mutation rates and the simulated-annealing
  schedule.}
  \label{fig:suite04-trap-box}
\end{figure}

Across Trap instances, \acrshort{sa} does not reach the global optimum within
the allotted budget, matching the failure of the \acrshort{ea} variants and
indicating that the tested temperature schedule is insufficient to
traverse the deceptive gradient (Figure~\ref{fig:suite04-trap-box}).
Its higher exploratory mobility
instead produces consistently larger incumbent cumulative regret: for
$k \neq 20$, the separation between the \acrshort{sa} trajectory and
the \acrshort{ea}
trajectories becomes slightly visible as $k$ grows, although the
eventual stabilized regret slopes remain broadly similar across
algorithms and deception widths.
The $d=8.0$ override on Trap-$k20$
does not alter this qualitative conclusion; it still fails to
converge within the budget and its \acrshort{cr} remains
substantially displaced from the \acrshort{ea} curves, again reflecting
exploration rather than successful escape. Thus, Suite~04 provides a
negative but informative result: increased exploration raises
cumulative regret without guaranteeing optimization success on the fully
deceptive Trap landscape (Figures~\ref{fig:suite04-trap-cr}
and~\ref{fig:suite04-trap-box}).

\subsection{Suites 05 and 09: Jump --- Mutation Rate and
Matched-Budget Reference}
\label{sec:results:suite05}
\begin{figure}[!htbp]
  \centering
  \begin{subfigure}[t]{0.48\textwidth}
    \includegraphics[width=\linewidth]{results/barrier_mut/jump/k2-history_regret_instantaneous_best.pdf}
    \caption{\textsc{Jump}-$2$.}
  \end{subfigure}
  \begin{subfigure}[t]{0.48\textwidth}
    \includegraphics[width=\linewidth]{results/barrier_mut/jump/k3-history_regret_instantaneous_best.pdf}
    \caption{\textsc{Jump}-$3$.}
  \end{subfigure}
  \caption[Suite 05 small-jump instantaneous regret]{Suite~05
    incumbent \acrshort{ir} trajectories for small
    jump widths. The plateau height identifies the residual barrier after the
  local optimum is reached.}
  \label{fig:jumpinst-small}
\end{figure}

\begin{figure}[!htbp]
  \centering
  \begin{subfigure}[t]{0.48\textwidth}
    \includegraphics[width=\linewidth]{results/barrier_mut/jump/k5-history_regret_instantaneous_best.pdf}
    \caption{\textsc{Jump}-$5$.}
  \end{subfigure}
  \begin{subfigure}[t]{0.48\textwidth}
    \includegraphics[width=\linewidth]{results/barrier_mut/jump/k7-history_regret_instantaneous_best.pdf}
    \caption{\textsc{Jump}-$7$.}
  \end{subfigure}
  \begin{subfigure}[t]{0.80\textwidth}
    \includegraphics[width=\linewidth]{results/barrier_mut/jump/k10-history_regret_instantaneous_best.pdf}
    \caption{\textsc{Jump}-$10$.}
  \end{subfigure}
  \caption[Suite 05 Jump instantaneous regret]{Suite~05 incumbent
    \acrshort{ir} trajectories for
    \textsc{Jump}-$k$, $k\in\{5,7,10\}$. The trajectories descend into a
    width-dependent plateau whose height equals the remaining fitness gap to
  the optimum.}
  \label{fig:jumpinst}
\end{figure}
Figures~\ref{fig:jumpinst-small} and~\ref{fig:jumpinst} depict the incumbent, instantaneous regret
curves for Jump-$k$, $k \in \{2,3,5,7,10\}$.
The instantaneous incumbent-regret trajectories show two
distinct phases: an initial descent toward the local optimum, followed by a
flat plateau at regret $k$ once the search reaches Hamming weight $n-k$.
On linear axes, increasing the mutation rate delays this settling phase:
higher-rate variants spend longer at intermediate incumbent-regret levels
before collapsing onto the plateau. Thus, the plateau height remains
structurally determined by the jump width, while the mutation rate primarily
changes the transient path into the plateau.

\begin{figure}[!htbp]
  \centering
  \begin{subfigure}[t]{0.48\textwidth}
    \includegraphics[width=\linewidth]{results/barrier_mut/jump/k2-history_regret_cumulative_best.pdf}
    \caption{\textsc{Jump}-$2$.}
    \label{fig:suite05-jump-k2-cr}
  \end{subfigure}
  \begin{subfigure}[t]{0.48\textwidth}
    \includegraphics[width=\linewidth]{results/barrier_mut/jump/k3-history_regret_cumulative_best.pdf}
    \caption{\textsc{Jump}-$3$.}
    \label{fig:suite05-jump-k3-cr}
  \end{subfigure}
  \begin{subfigure}[t]{0.48\textwidth}
    \includegraphics[width=\linewidth]{results/barrier_mut/jump/k5-history_regret_cumulative_best.pdf}
    \caption{\textsc{Jump}-$5$.}
    \label{fig:suite05-jump-k5-cr}
  \end{subfigure}
  \begin{subfigure}[t]{0.48\textwidth}
    \includegraphics[width=\linewidth]{results/barrier_mut/jump/k7-history_regret_cumulative_best.pdf}
    \caption{\textsc{Jump}-$7$.}
    \label{fig:suite05-jump-k7-cr}
  \end{subfigure}
  \begin{subfigure}[t]{0.80\textwidth}
    \includegraphics[width=\linewidth]{results/barrier_mut/jump/k10-history_regret_cumulative_best.pdf}
    \caption{\textsc{Jump}-$10$.}
    \label{fig:suite05-jump-k10-cr}
  \end{subfigure}
  \caption[Suite 05 Jump cumulative regret]{Suite~05 incumbent
    \acrshort{cr} trajectories on
    \textsc{Jump} as the jump width increases. After the initial descent, the
    trapped phase appears as an approximately linear ramp whose slope is set by
  the incumbent regret plateau.}
  \label{fig:suite05-jump-cr}
\end{figure}

\begin{figure}[!htbp]
  \centering
  \begin{subfigure}[t]{0.48\textwidth}
    \includegraphics[width=\linewidth]{results/barrier_mut/jump_plateau/k2-history_regret_cumulative_best.pdf}
    \caption{Matched-budget cumulative regret.}
    \label{fig:suite09-jump-k2-cr}
  \end{subfigure}
  \begin{subfigure}[t]{0.48\textwidth}
    \includegraphics[width=\linewidth]{results/barrier_mut/jump_plateau/k2-convergence_probability.pdf}
    \caption{Matched-budget convergence probability.}
    \label{fig:suite09-jump-k2-cdp}
  \end{subfigure}
  \caption[Suite 09 matched-budget Jump reference]{Suite~09
    matched-budget reference for \textsc{Jump}-$2$.
    Cumulative regret records the cost of the trapped phase, while convergence
  probability records only whether the optimum has been hit by each budget.}
  \label{fig:suite09-jump-k2}
\end{figure}
The adaptive $k/n$ mutation rate increasingly separates from the fixed-rate
\acrshort{cr} curves as $k$ grows, a pattern visible both in Suite~05 and
in the matched-budget Suite~09 reference. In Suite~05, the $k/n$ variant
coincides with the fixed $2/n$ rate at $k=2$; for larger $k$, the longer
residence at intermediate incumbent-regret levels before the regret-$k$ plateau
produces a larger \acrshort{cr} gap; see
Figures~\ref{fig:suite05-jump-cr} and~\ref{fig:suite09-jump-k2}. This
also shows why mean simple regret and convergence probability can be
misleading at high budgets: they record the eventual plateau or success state,
but not the accumulated cost of reaching it.

\subsection{Suite 06: Plateau --- Neutral Walk and Mutation Rate}
\label{sec:results:suite06}
\begin{figure}[!htbp]
  \centering
  \begin{subfigure}[t]{0.48\textwidth}
    \includegraphics[width=\linewidth]{results/barrier_mut/plateau/k2-history_regret_instantaneous_best.pdf}
    \caption{\textsc{Plateau}-$2$ instantaneous regret.}
    \label{fig:suite06-plateau-k2-ir}
  \end{subfigure}
  \begin{subfigure}[t]{0.48\textwidth}
    \includegraphics[width=\linewidth]{results/barrier_mut/plateau/k2-history_regret_cumulative_best.pdf}
    \caption{\textsc{Plateau}-$2$ cumulative regret.}
    \label{fig:suite06-plateau-k2-cr}
  \end{subfigure}
  \begin{subfigure}[t]{0.48\textwidth}
    \includegraphics[width=\linewidth]{results/barrier_mut/jump_plateau/k2-history_regret_cumulative_best.pdf}
    \caption{\textsc{Jump}-$2$ cumulative regret.}
    \label{fig:suite09-jump-k2-cr-plateau-ref}
  \end{subfigure}
  \begin{subfigure}[t]{0.48\textwidth}
    \includegraphics[width=\linewidth]{results/barrier_mut/jump_plateau/k2-convergence_probability.pdf}
    \caption{\textsc{Jump}-$2$ convergence probability.}
    \label{fig:suite09-jump-k2-cdp-plateau-ref}
  \end{subfigure}
  \caption[Matched Plateau and Jump comparison]{Matched $k=2$
    comparison between the neutral-walk
    \textsc{Plateau} landscape and the fitness-valley \textsc{Jump}
    landscape. Plateau and Jump differ most clearly in the early
  \acrshort{ir} descent and in the median \acrshort{cr} recovery.}
  \label{fig:suite06-plateau-jump-k2}
\end{figure}

\begin{figure}[!htbp]
  \centering
  \includegraphics[width=0.80\textwidth]{results/barrier_mut/plateau_jump/plateau_vs_jump_k5_b10000.pdf}
  \caption[Plateau and Jump regret distributions]{Budget-$10{,}000$
    comparison of \textsc{Plateau}-$5$ and
    \textsc{Jump}-$5$ regret distributions. At this width the neutral-walk and
    fitness-valley mechanisms are much less visually separable than at
  $k=2$.}
  \label{fig:suite06-plateau-jump-k5}
\end{figure}

The Plateau results are most informative when compared cross-suite against the
matched-budget Jump reference in Suite~09, since Suite~06 does not
contain Jump curves.
At $k=2$, the incumbent \acrshort{cr} trajectories are visually distinct:
Plateau-$2$ shows faster median recovery for several EA variants, and under the
$10{,}000$-evaluation budget some median curves reach zero incumbent
instantaneous regret (Figure~\ref{fig:suite06-plateau-k2-ir}),
indicating convergence to the global optimum
in the \acrshort{cr} trajectory
(Figure~\ref{fig:suite06-plateau-k2-cr}). This contrasts with
Jump-$2$ at the same budget, where optima are
sampled in some runs, but the median trajectory does not converge
Figures~\ref{fig:suite09-jump-k2-cr-plateau-ref}
and~\ref{fig:suite09-jump-k2-cdp-plateau-ref}. At $k=5$, however, the
distinction largely disappears:
the Plateau and Jump \acrshort{cr} curves are similar both in shape and in
magnitude (Figure~\ref{fig:suite06-plateau-jump-k5}). Thus, the
neutral-walk mechanism is
clearly separable from Jump at small $k$, but becomes empirically harder to
distinguish from a fitness-valley barrier as $k$ increases.

Thus, Suite~06 exposes a boundary of the \acrshort{cr} diagnostic:
\acrshort{cr} is not meant to distinguish the landscape identity, it just
hints at the realized optimization cost.
When Plateau and Jump differ strongly, as at $k=2$, their
incumbent \acrshort{cr} trajectories separate; when their induced search
costs are similar, as at $k=5$, the curves may overlap despite different
underlying mechanisms. Cumulative regret should therefore be read as an
algorithm--problem trajectory diagnostic.

\subsection{Suite 10: Barrier Archetype Comparison at $k=5$}
\label{sec:results:suite10}

Suite~10 fixes $k=5$ across \textsc{Jump}, \textsc{Trap}, and
\textsc{Plateau} because endpoint convergence is deliberately absent:
none of the tested $(1+1)$-\acrshort{ea} mutation-rate variants reaches
the global optimum within the available budget on any of the three
problems. Rather than reproduce three identical zero-convergence panels,
Figure~\ref{fig:suite10-regret-boxplot} shows the corresponding budget-level
regret distribution for the representative \textsc{Jump}-$5$ case.
The comparison therefore queries how different barrier mechanisms are
reflected in the regret trajectories when \acrshort{ttfo} provides no
usable samples.

\begin{figure}[!htbp]
  \centering
  \includegraphics[width=0.80\textwidth]{results/barrier_mut/all_prob/jump_k5_regret_boxplot_b10000.pdf}
  \caption[Suite 10 representative regret distribution]{Suite~10
    budget-$10{,}000$ regret distribution for the representative
    \textsc{Jump}-$5$ archetype. The omitted convergence-probability panels for
    \textsc{Jump}-$5$, \textsc{Plateau}-$5$, and \textsc{Trap}-$5$ are exactly
    identical for the purpose used here: they show zero endpoint convergence for
    every mutation-rate variant, so the regret distribution is the more
    informative finite-budget diagnostic.}
  \label{fig:suite10-regret-boxplot}
\end{figure}

\begin{figure}[!htbp]
  \centering
  \begin{subfigure}[t]{0.48\textwidth}
    \includegraphics[width=\linewidth]{results/barrier_mut/all_prob/jump_k5_history_regret_instantaneous_best.pdf}
    \caption{\textsc{Jump}-$5$ instantaneous regret.}
    \label{fig:suite10-jump-k5-ir}
  \end{subfigure}
  \begin{subfigure}[t]{0.48\textwidth}
    \includegraphics[width=\linewidth]{results/barrier_mut/all_prob/plateau_k5_history_regret_instantaneous_best.pdf}
    \caption{\textsc{Plateau}-$5$ instantaneous regret.}
    \label{fig:suite10-plateau-k5-ir}
  \end{subfigure}
  \begin{subfigure}[t]{0.80\textwidth}
    \includegraphics[width=\linewidth]{results/barrier_mut/all_prob/trap_k5_history_regret_instantaneous_best.pdf}
    \caption{\textsc{Trap}-$5$ instantaneous regret.}
    \label{fig:suite10-trap-k5-ir}
  \end{subfigure}
  \caption[Suite 10 barrier-archetype instantaneous regret]{Suite~10
    incumbent \acrshort{ir} trajectories for the three $k=5$ barrier
    archetypes. The plots expose the trapped-regret floors even when
    convergence probability is uniformly uninformative.}
  \label{fig:suite10-ir}
\end{figure}

\begin{figure}[!htbp]
  \centering
  \begin{subfigure}[t]{0.48\textwidth}
    \includegraphics[width=\linewidth]{results/barrier_mut/all_prob/jump_k5_history_regret_cumulative_best.pdf}
    \caption{\textsc{Jump}-$5$ cumulative regret.}
    \label{fig:suite10-jump-k5-cr}
  \end{subfigure}
  \begin{subfigure}[t]{0.48\textwidth}
    \includegraphics[width=\linewidth]{results/barrier_mut/all_prob/plateau_k5_history_regret_cumulative_best.pdf}
    \caption{\textsc{Plateau}-$5$ cumulative regret.}
    \label{fig:suite10-plateau-k5-cr}
  \end{subfigure}
  \begin{subfigure}[t]{0.80\textwidth}
    \includegraphics[width=\linewidth]{results/barrier_mut/all_prob/trap_k5_history_regret_cumulative_best.pdf}
    \caption{\textsc{Trap}-$5$ cumulative regret.}
    \label{fig:suite10-trap-k5-cr}
  \end{subfigure}
  \caption[Suite 10 barrier-archetype cumulative regret]{Suite~10
    incumbent \acrshort{cr} trajectories for the three $k=5$ barrier
    archetypes. The cumulative curves record the cost of each trapped phase
    under identical mutation-rate variants and budget.}
  \label{fig:suite10-cr}
\end{figure}

The \textsc{Jump}-$5$ and \textsc{Plateau}-$5$ incumbent
\acrshort{cr} curves remain close in shape, matching the behavior
already observed in the Suite~06 matched comparison~\ref{fig:suite06-plateau-jump-k2}. At this
width, the neutral-walk and gap-crossing mechanisms induce similar
realized \acrshort{ir} curves under the tested mutation rates, so their cumulative
curves primarily record comparable optimization cost rather than a
visibly separate barrier signature (Figures~\ref{fig:suite10-ir}
and~\ref{fig:suite10-cr}).

The \textsc{Trap}-$5$ median incumbent \acrshort{ir} trajectories
show a clearer contrast. Across all four mutation-rate variants, the
algorithm spends substantially longer at intermediate fitness
levels before settling into the deceptive basin with residual regret
approximately $k$ (Figure~\ref{fig:suite10-trap-k5-ir}). This extended intermediate phase is
consistent with the deceptive gradient: progress can continue locally,
but it carries the search toward the wrong basin rather than toward the
global optimum.

Thus, although no method solves any of the three instances within the
budget, the trajectory plots still expose how each barrier shapes the
cost of failure.

\section{NK Ruggedness and SA Cooling}
\label{sec:results:nk}

\subsection{Suite 11: SA Cooling Schedule Diagnostics}
\label{sec:results:suite11}
\begin{figure}[!htbp]
  \centering
  \includegraphics[width=0.80\textwidth]{results/nk_sa/sa_cooling_onemax_history_regret_instantaneous.pdf}
  \caption[Suite 11 SA cooling trajectories]{Suite~11 incumbent
    \acrshort{ir} trajectories on
    \textsc{OneMax} for \acrshort{rls} and simulated-annealing cooling
    schedules. The onset of each low-regret floor gives a direct visual
  diagnostic of effective cooling speed.}
  \label{fig:suite11-sa-cooling-ir}
\end{figure}

The instantaneous regret trajectories show a clear ordering in the
effective cooling, or for RLS, stabilization, time of the methods.
The empirical ordering is
\begin{multline*}
  \text{RLS} \;<\; \text{SA-Exp} \;<\; \text{SA-Log}_{d=0.5} \;<\;
  \text{SA-Log}_{d=1} \\ \;<\; \text{SA-Log}_{d=2} \;<\;
  \text{SA-Log}_{d=4} \;<\; \text{SA-Lin} \;<\; \text{SA-Log}_{d=8}
  \;<\; \text{SA-Log}_{d=16},
\end{multline*}
where earlier position denotes earlier convergence to a low-regret
plateau. Thus, the log-scale \acrshort{ir} curves provide a direct
diagnostic of schedule aggressiveness: fast schedules collapse
quickly, while slower logarithmic schedules preserve exploration
longer and delay stabilization. This supports the suite hypothesis
that cooling schedules leave visible trajectory-level signatures,
rather than only changing final regret values; see
Figure~\ref{fig:suite11-sa-cooling-ir}.

\subsection{Suite 03: NK Ruggedness --- Cooling $\times$ Epistasis Grid}
\label{sec:results:suite03}

\begin{figure}[!htbp]
  \centering
  \begin{subfigure}[t]{0.48\textwidth}
    \includegraphics[width=\linewidth]{results/nk_sa/suite03_k1_history_regret_cumulative_best.pdf}
    \caption{NK landscape with $k=1$.}
    \label{fig:suite03-k1-cr}
  \end{subfigure}
  \begin{subfigure}[t]{0.48\textwidth}
    \includegraphics[width=\linewidth]{results/nk_sa/suite03_k2_history_regret_cumulative_best.pdf}
    \caption{NK landscape with $k=2$.}
    \label{fig:suite03-k2-cr}
  \end{subfigure}
  \begin{subfigure}[t]{0.80\textwidth}
    \includegraphics[width=\linewidth]{results/nk_sa/suite03_k4_history_regret_cumulative_best.pdf}
    \caption{NK landscape with $k=4$.}
    \label{fig:suite03-k4-cr}
  \end{subfigure}
  \caption[Suite 03 NK \acrshort{cr} trajectories]{Suite~03
    incumbent \acrshort{cr} trajectories across
    increasing NK epistasis. Low-epistasis instances favour mutation-only
    baselines, while at $k=4$ the slower logarithmic cooling schedule begins to
  dominate after a longer transient phase.}
  \label{fig:suite03-cr}
\end{figure}

\begin{figure}[!htbp]
  \centering
  \begin{subfigure}[t]{0.48\textwidth}
    \includegraphics[width=\linewidth]{results/nk_sa/suite03_k4_inverse_runtime_profile_surface_RLS.pdf}
    \caption{\acrshort{rls} on NK $k=4$.}
    \label{fig:suite03-k4-profile-rls}
  \end{subfigure}
  \begin{subfigure}[t]{0.48\textwidth}
    \includegraphics[width=\linewidth]{results/nk_sa/suite03_k4_inverse_runtime_profile_surface_SA-Log-d_0.5.pdf}
    \caption{\acrshort{sa}-Log with $d=0.5$ on NK $k=4$.}
    \label{fig:suite03-k4-profile-sa-log-d05}
  \end{subfigure}
  \caption[Suite 03 NK inverse runtime profiles]{Representative
    Suite~03 inverse runtime profile surfaces for NK
    $k=4$. The saturation frontiers show where each algorithm ceases to make
  reliable progress toward higher fitness levels.}
  \label{fig:suite03-profiles}
\end{figure}

Suite~03 shows a clear regime split between low and moderate
epistasis. For $k=1$ and $k=2$, all \acrshort{sa} variants underperform the
baselines in incumbent cumulative regret, indicating that
cooling-induced exploration is not compensated by ruggedness at these
settings (Figure~\ref{fig:suite03-k2-cr}). The effect is strongest at
$k=1$, where the $(1+1)$-\acrshort{ea} is
the only method whose median incumbent \acrshort{cr} curve
reaches a horizontal asymptote at the optimum, stabilizing at
approximately $15.63$ cumulative regret (Figure~\ref{fig:suite03-k1-cr}).
This ordering reverses for $k \ge 4$: \acrshort{sa} with
logarithmic cooling at $d=0.5$ accumulates the least regret overall
and separates clearly from RLS and the $(1+1)$-\acrshort{ea}, although no median
curve converges within the evaluation budget, as expected on these
rugged instances. At $k=4$, the \acrshort{sa}-\acrshort{log}-$d=0.5$
curve crosses below
the $(1+1)$-\acrshort{ea} at roughly $3.3\times 10^4$ evaluations, marking the
budget at which slower exploratory progress begins to dominate greedy
improvement (Figure~\ref{fig:suite03-k4-cr}). The inverse runtime profile
surfaces corroborate
this interpretation: across the grid, \acrshort{rls} exhibits a truncated
frontier consistent with entrapment in a local optimum, visible from
the contour boundary of the corresponding surface
(Figure~\ref{fig:suite03-profiles}).
% ~200 words
%
\subsection{Suite 12: NK Saturation}
\label{sec:results:suite12}
\begin{figure}[!htbp]
  \centering
  \begin{subfigure}[t]{0.48\textwidth}
    \includegraphics[width=\linewidth]{results/nk_sa/suite12_k0_history_regret_instantaneous_best.pdf}
    \caption{NK $k=0$.}
    \label{fig:suite12-ir-k0}
  \end{subfigure}
  \begin{subfigure}[t]{0.48\textwidth}
    \includegraphics[width=\linewidth]{results/nk_sa/suite12_k2_history_regret_instantaneous_best.pdf}
    \caption{NK $k=2$.}
    \label{fig:suite12-ir-k2}
  \end{subfigure}
  \begin{subfigure}[t]{0.48\textwidth}
    \includegraphics[width=\linewidth]{results/nk_sa/suite12_k4_history_regret_instantaneous_best.pdf}
    \caption{NK $k=4$.}
    \label{fig:suite12-ir-k4}
  \end{subfigure}
  \begin{subfigure}[t]{0.48\textwidth}
    \includegraphics[width=\linewidth]{results/nk_sa/suite12_k6_history_regret_instantaneous_best.pdf}
    \caption{NK $k=6$.}
    \label{fig:suite12-ir-k6}
  \end{subfigure}
  \begin{subfigure}[t]{0.80\textwidth}
    \includegraphics[width=\linewidth]{results/nk_sa/suite12_k8_history_regret_instantaneous_best.pdf}
    \caption{NK $k=8$.}
    \label{fig:suite12-ir-k8}
  \end{subfigure}
  \caption[Suite 12 calibrated-cooling regret trajectories]{Suite~12
    incumbent \acrshort{ir} trajectories under
    per-$k$ calibrated cooling. The terminal regret floor rises with epistasis,
    and calibrated simulated annealing becomes most useful on the more rugged
  $k=6$ and $k=8$ landscapes.}
  \label{fig:suite12-ir}
\end{figure}

\begin{figure}[!htbp]
  \centering
  \begin{subfigure}[t]{0.48\textwidth}
    \includegraphics[width=\linewidth]{results/nk_sa/suite12_k6_inverse_runtime_profile_surface_RLS.pdf}
    \caption{\acrshort{rls} on NK $k=6$.}
    \label{fig:suite12-k6-profile-rls}
  \end{subfigure}
  \begin{subfigure}[t]{0.48\textwidth}
    \includegraphics[width=\linewidth]{results/nk_sa/suite12_k6_inverse_runtime_profile_surface_SA-Log-d_8.0.pdf}
    \caption{Calibrated \acrshort{sa} on NK $k=6$.}
    \label{fig:suite12-k6-profile-sa}
  \end{subfigure}
  \begin{subfigure}[t]{0.48\textwidth}
    \includegraphics[width=\linewidth]{results/nk_sa/suite12_k8_inverse_runtime_profile_surface_RLS.pdf}
    \caption{\acrshort{rls} on NK $k=8$.}
    \label{fig:suite12-k8-profile-rls}
  \end{subfigure}
  \begin{subfigure}[t]{0.48\textwidth}
    \includegraphics[width=\linewidth]{results/nk_sa/suite12_k8_inverse_runtime_profile_surface_SA-Log-d_12.0.pdf}
    \caption{Calibrated \acrshort{sa} on NK $k=8$.}
    \label{fig:suite12-k8-profile-sa}
  \end{subfigure}
  \caption[Suite 12 high-epistasis NK profiles]{Suite~12 inverse
    runtime profile surfaces for higher-epistasis NK
    landscapes. The RLS surfaces saturate below the optimum, while calibrated
    simulated annealing maintains progress to higher fitness levels under the
  displayed budgets.}
  \label{fig:suite12-profiles}
\end{figure}

In Suite~12 (Table~\ref{tab:suite12-d}), the incumbent
\acrshort{ir} trajectories show a
clear regime change with increasing epistasis. For $k \in \{0,2,4\}$,
\acrshort{sa} does not fully collapse within the allotted budget
(exploring), and the
$(1+1)$-\acrshort{ea} eventually attains lower regret on several instances; this
is visible in the corresponding best-history curves
Figures~\ref{fig:suite12-ir-k0}--\ref{fig:suite12-ir-k4}. At higher
ruggedness the ordering reverses: for $k=6$, calibrated SA with $d=8$
moves below the EA trajectory around evaluation $\approx 8500$, while
for $k=8$, calibrated SA with $d=12$ retains a lower terminal regret
at budget $25000$ (Figures~\ref{fig:suite12-ir-k6}
and~\ref{fig:suite12-ir-k8}). Thus, the calibrated cooling schedule
is not uniformly advantageous on weakly rugged landscapes, but
becomes useful once mutation-only search is trapped in more
suboptimal local optima. The inverse runtime profile surfaces
corroborate this interpretation: in particular, the RLS frontier
saturates below the optimum, and the contour boundary gives a direct
visual indication of the local fitness level at which the algorithm
becomes trapped (Figures~\ref{fig:suite12-k6-profile-rls},
  \ref{fig:suite12-k6-profile-sa}, \ref{fig:suite12-k8-profile-rls},
and~\ref{fig:suite12-k8-profile-sa}).
\section{Hierarchical Structure}
\label{sec:results:hierarchical}

\subsection{Suite 07: \texorpdfstring{\acrshort{hiff}}{HIFF} Scaling
with Problem Size}
\label{sec:results:suite07}

\begin{figure}[!htbp]
  \centering
  \begin{subfigure}[t]{0.48\textwidth}
    \includegraphics[width=\linewidth]{results/hiff/n32_history_regret_instantaneous_best.pdf}
    \caption{$n=32$ instantaneous regret.}
    \label{fig:suite07-n32-ir}
  \end{subfigure}
  \begin{subfigure}[t]{0.48\textwidth}
    \includegraphics[width=\linewidth]{results/hiff/n64_history_regret_instantaneous_best.pdf}
    \caption{$n=64$ instantaneous regret.}
    \label{fig:suite07-n64-ir}
  \end{subfigure}
  \begin{subfigure}[t]{0.80\textwidth}
    \includegraphics[width=\linewidth]{results/hiff/n128_history_regret_instantaneous_best.pdf}
    \caption{$n=128$ instantaneous regret.}
    \label{fig:suite07-n128-ir}
  \end{subfigure}
  \caption[Suite 07 HIFF instantaneous regret]{Suite~07 incumbent
    \acrshort{ir} trajectories for \acrshort{hiff} as problem size increases.
    The terminal floors show the partial-progress level reached under the
    finite evaluation budgets.}
  \label{fig:suite07-ir}
\end{figure}

\begin{figure}[!htbp]
  \centering
  \begin{subfigure}[t]{0.48\textwidth}
    \includegraphics[width=\linewidth]{results/hiff/n32_history_best.pdf}
    \caption{$n=32$ best-so-far fitness.}
    \label{fig:suite07-n32-best}
  \end{subfigure}
  \begin{subfigure}[t]{0.48\textwidth}
    \includegraphics[width=\linewidth]{results/hiff/n64_history_best.pdf}
    \caption{$n=64$ best-so-far fitness.}
    \label{fig:suite07-n64-best}
  \end{subfigure}
  \begin{subfigure}[t]{0.80\textwidth}
    \includegraphics[width=\linewidth]{results/hiff/n128_history_best.pdf}
    \caption{$n=128$ best-so-far fitness.}
    \label{fig:suite07-n128-best}
  \end{subfigure}
  \caption[Suite 07 HIFF best-history plots]{Suite~07 best-so-far fitness
    histories for \acrshort{hiff}. These plots show the hierarchical
    stagnation levels directly, complementing the regret view.}
  \label{fig:suite07-best}
\end{figure}

\begin{figure}[!htbp]
  \centering
  \begin{subfigure}[t]{0.48\textwidth}
    \includegraphics[width=\linewidth]{results/hiff/n32_history_regret_cumulative_best.pdf}
    \caption{$n=32$ cumulative regret.}
    \label{fig:suite07-n32-cr}
  \end{subfigure}
  \begin{subfigure}[t]{0.48\textwidth}
    \includegraphics[width=\linewidth]{results/hiff/n64_history_regret_cumulative_best.pdf}
    \caption{$n=64$ cumulative regret.}
    \label{fig:suite07-n64-cr}
  \end{subfigure}
  \begin{subfigure}[t]{0.80\textwidth}
    \includegraphics[width=\linewidth]{results/hiff/n128_history_regret_cumulative_best.pdf}
    \caption{$n=128$ cumulative regret.}
    \label{fig:suite07-n128-cr}
  \end{subfigure}
  \caption[Suite 07 HIFF cumulative regret]{Suite~07 incumbent
    \acrshort{cr} trajectories for \acrshort{hiff}. Increasing problem size
    raises the cumulative cost of partial progress under the fixed budgets.}
  \label{fig:suite07-cr}
\end{figure}

Suite~07 should be read as a finite-budget stress test of hierarchical
structure, not as a dedicated asymptotic scaling experiment. Across the tested
problem sizes, $n \in \{32,64,128\}$, the qualitative failure mode is strict
one-bit greediness: once \acrshort{rls} reaches a locally coherent block
structure, improving moves typically require coordinated flips that are
inaccessible to its neighborhood. The resulting \acrshort{ir} floors and
best-fitness stagnation make this failure mode visible before the final budget
(Figures~\ref{fig:suite07-ir} and~\ref{fig:suite07-best}).

The exploratory \acrshort{rlse} variant and the $(1+1)$-\acrshort{ea}
mutation-rate
grid test two ways of relaxing this limitation: accepting occasional worsening
moves, and allowing multi-bit flips. Over the three tested
\acrshort{hiff} sizes,
these mechanisms improve partial progress but do not establish a general
convergence threshold (Figures~\ref{fig:suite07-ir}
and~\ref{fig:suite07-best}). The appropriate conclusion is therefore narrower than a
regret-scaling law: on hierarchical landscapes, mutation and exploration affect
both the final incumbent quality and the rate at which cumulative regret accrues
under a fixed evaluation budget (Figures~\ref{fig:suite07-ir},
\ref{fig:suite07-best}, and~\ref{fig:suite07-cr}).

\section{Cross-Suite Synthesis}
\label{sec:results:synthesis}
%   is never reached within budget --- the condition under which classical
%   measures are least useful and anytime performance measurement matters most.

Across the experimental suites, three takeaways recur. They do not
show that \acrshort{cr} uniquely identifies a landscape class; rather,
they show that the best-so-far regret trajectory records the realized
cost of search under a finite evaluation budget.

\paragraph{Regret discriminates when \acrshort{ttfo} cannot.}
In the hard finite-budget cases, binary convergence measures often go
silent: \textsc{Trap} instances remain non-convergent within budget,
\textsc{Jump} instances with larger $k$ settle onto a width-dependent
regret plateau, and \textsc{Plateau} becomes difficult to distinguish
from \textsc{Jump} once the induced search cost is similar
(Figures~\ref{fig:suite04-trap-cr}, \ref{fig:suite05-jump-cr},
\ref{fig:suite06-plateau-jump-k2}, and~\ref{fig:suite06-plateau-jump-k5}).
High-epistasis NK and larger \acrshort{hiff} instances provide analogous
rugged and hierarchical failure modes rather than literal barrier
landscapes. In all three settings, the slope of incumbent \acrshort{cr},
the floor of incumbent \acrshort{ir}, and the onset of stagnation
continue to describe how far the retained solution remains from the
optimum. This is the operational content of the project's central claim:
regret provides structured information precisely where \acrshort{ttfo}
distributions and $P(\mathrm{opt})$ summaries are least informative.

\paragraph{The asymptotic regret ordering can invert the \acrshort{ttfo} ordering.}
Suites~01 and~02 show that earlier arrival at the global optimum does
not necessarily imply lower cumulative cost. The inversion on
\textsc{LeadingOnes} and \textsc{BinVal}, where \acrshort{rls}
reaches the optimum before the $(1+1)$-\acrshort{ea} but accumulates a
larger regret asymptote, follows the mechanism of
Corollary~\ref{cor:asymptotic-regret}. The asymptote is governed by
first-hitting times across \emph{all} intermediate fitness levels, not
only by the terminal hitting time. Thus, \acrshort{ttfo} records final
arrival, while incumbent cumulative regret records the area under the
best-so-far fitness gap: two algorithms can have the same endpoint, or
even opposite endpoint rankings, while paying different path costs.

\paragraph{Cooling schedule aggressiveness leaves a visible trajectory
signature.}
Suite~11 demonstrates that the \acrshort{ir} curve gives a direct
diagnostic of cooling speed, with the onset of the low-regret floor
ordering schedules from fastest (SA-Exp) to slowest (SA-Log-$d=16$)
(Figure~\ref{fig:suite11-sa-cooling-ir}). Suites~03 and~12 show why this
matters: on low-epistasis NK landscapes, temperature-driven acceptance
inflates cumulative regret relative to hill-climbing baselines, whereas
at high epistasis ($k \geq 6$), calibrated logarithmic cooling maintains
progress after mutation-only methods saturate
(Figures~\ref{fig:suite03-cr} and~\ref{fig:suite12-ir}). The budget at
which SA overtakes the $(1+1)$-\acrshort{ea} on NK-$k=4$ (roughly
$3.3 \times 10^4$ evaluations) is therefore best read as an analog of
the crossover behavior motivating runtime-profile reporting in
\citet{pinto2018practice}, expressed here in the \acrshort{cr} plane
rather than as a literal crossover fitness level.

Suite~08 supports the theoretical interpretation behind these empirical
readings: on integer-valued benchmarks, the profile-derived and directly
averaged \acrshort{cr} curves agree to floating-point precision, consistent with
the tail-sum decomposition of Theorem~\ref{thm:regret-profile}. For real-valued
landscapes such as NK, the profile surfaces should still be read as discretized
approximations, a limitation that bounds but does not invalidate the qualitative
conclusions drawn in Section~\ref{sec:results:nk}. The practical conclusion is that
\acrshort{cr} should be read as a trajectory diagnostic: its slope,
floor, and crossover points expose failure modes that \acrshort{ttfo}
and convergence probability hide under finite budgets.
